\subsection*{15. Test Potential by Teaching Something, Not Asking Whether Someone Likes Learning}

``Are you a fast learner?'' is an invitation to answer yes. A stronger method is a \textbf{dynamic assessment}: observe the candidate's initial approach to a novel, job-relevant problem; supply a standard hint, missing fact, or piece of feedback; then offer a parallel problem and ask them to explain what changed. Score baseline performance, quality of questions, correction speed, generalization, and ability to state remaining uncertainty separately.

For instance, give 2 engineers an unfamiliar logging format. Both miss the 1st failure mode. Candidate A, after a hint about clock skew, checks timestamps, forms a new hypothesis, tests it, and applies the idea to a 2nd trace. Candidate B memorizes the hint and repeats it without inspecting the 2nd trace. The difference is \textbf{observed learning behavior}, not a speculative ``potential score.'' A single demonstration is still noisy; use it as a bounded observation, especially for an apprentice role.

\subsection*{16. Consider the Difference Between Learning Rate and Learning Opportunity}

Candidates arrive with unequal access to prestigious employers, computing resources, mentors, professional networks, English fluency, or expensive training. Those differences can affect visible polish without reflecting their ultimate capacity. An employer need not ignore job requirements; it should distinguish what was \textbf{learned through opportunity} from what the person is likely to learn with the opportunity the organization can actually provide.

Ask for examples of building capability under resource constraints. Give all finalists equivalent instructions, a chance to ask clarifying questions, and accessible assessment conditions. Avoid interpreting a candidate's unfamiliarity with your preferred tool as an absence of systems thinking. Equally, do not romanticize ``raw potential'' when the role has no supervisor, no learning budget, and an immediate critical deadline. Hiring potential without the ability to develop it is an unfair promise to both sides.

\subsection*{17. Use Must-Pass Gates Before Weighted Averages}

A candidate who scores 5/5 in creativity but 1/5 in data privacy should not become a customer-facing AI engineer because the total average is high. Identify non-compensatory gates for legal compliance, safety, honesty, domain competence, or vulnerability of customers. Only after candidates clear gates should weighted strengths decide among them.

For example, technical reviewer competencies might be:
\begin{itemize}[leftmargin=*]
\item Technical correctness: 30\%
\item Diagnostic depth: 25\%
\item Actionable feedback: 20\%
\item Turnaround reliability: 15\%
\item Collaboration: 10\%
\end{itemize}
with a \textbf{minimum gate of 3/5 on correctness and integrity}. These weights are illustrative, not calibrated predictions; they should be revised with evidence. It is defensible to prefer fewer weights when differences cannot be justified; weighting becomes dangerous when the formula creates false precision.

\subsection*{18. A Scorecard Is a Record of Evidence, Not a Laundering Machine for Intuition}

For every score, record a verbatim behavior or artifact: ``missed the false-positive leakage in case B even after reading evaluation notes''; ``asked 3 questions about learner prerequisites before teaching''; ``explained to the buyer why the requested guarantee could not be made.'' Record confidence (high, medium, or low) and dissenting observations. A high score supported only by ``seemed strong'' should be treated as missing evidence.

Separately record \textbf{candidate interest, availability, compensation expectations, work authorization, conflicts, and references}. Fit of terms is not the same as skill. A technically excellent person who cannot work the required hours might be an excellent person and an infeasible hire. Keep those judgments apart.

\subsection*{19. Do Not Substitute a Prestigious Proxy for a Target Outcome}

Degrees, institutional brands, prior employers, publications, technical certificates, and revenue histories can indicate exposure to relevant work. They are rarely sufficient proof of the specific behavior required. A degree says something about prior selection and education; it does not directly prove the candidate can teach a difficult concept, debug a production pipeline, or close a complex procurement deal ethically.

Use credentials as \textbf{entry signals when justified}, then test the work. Some fields legitimately require licenses or degrees. For other fields, preferentially examine demonstrable competence and learning. Never claim that portfolios or credential-blind hiring are universally superior: the right policy depends on job analysis, validity evidence, legal conditions, and candidate access.

\subsection*{20. Standardize the Important Things, Not the Performative Ones}

Every candidate should receive comparable task information, deadlines, evaluation criteria, and reasonable accommodations. But forcing every specialist to answer the same generic question would be misguided. A researcher needs an epistemic exercise; a salesperson needs an interactive buyer situation; a teacher needs a learning encounter. \textbf{Fairness within a role} is compatible with \textbf{different methods across roles}.

For high-judgment tasks, standardize the scenario and scoring while permitting natural problem-solving approaches. Observe how the person frames the problem, requests missing facts, chooses between methods, checks work, and communicates limitations. This preserves useful variation rather than rewarding mere memorization of the hiring team's preferred script.
